% Pile de concentration au cuivre (exercice 14, chapitre 8)
\begin{tikzpicture}
\pcPile{pcCuivre}{pcBleu!30}{pcCuivre}{pcBleu}
\draw[pcFil] (pileG) -- (-0.4,4.4) -- (6.6,4.4) -- (pileD);
\pcAppareilRond{(3.1,4.4)}{V}
\node[pcLegende,anchor=south east] at (2.72,4.45) {COM};
\node[pcLegende,anchor=south west] at (3.48,4.45) {V};
\node[pcLegendeG,anchor=south] at (3.1,4.85) {tension affichée : $+\qty{0,030}{V}$};
\node[pcLegende] at (3.1,2.3) {pont salin};
\node[pcLegende,anchor=north,align=center] at (0,-0.15) {compartiment 1\\$[\ce{Cu^2+}]=\qty{1,0e-2}{mol.L^{-1}}$};
\node[pcLegende,anchor=north,align=center] at (6.2,-0.15) {compartiment 2\\$[\ce{Cu^2+}]=\qty{0,10}{mol.L^{-1}}$};
\end{tikzpicture}
\par\smallskip\legende{Deux demi-piles au cuivre qui ne diffèrent que par la concentration de leur solution.}
